\section{Form: Mengambil Nilai Input dan Validasi di Client}

Validasi form di client-side memberikan umpan balik instan tanpa menunggu respons server. Nilai input diakses melalui \code{.value} (teks, email, password), \code{.checked} (checkbox, radio), dan \code{.selectedIndex}/\code{.value} (select). Validasi dapat dilakukan saat event \code{input} (real-time), \code{change} (setelah blur), atau \code{submit} \cite{jsInfo}.

\begin{contoh}
\textbf{Studi Kasus: Form Registrasi dengan Validasi Real-Time}

Form dengan tiga field: email, password, dan konfirmasi password. Validasi berjalan real-time saat pengguna mengetik.
\begin{enumerate}
  \item Pasang \code{addEventListener("input", validate)} di setiap field
  \item Validasi email: harus mengandung ``@'' dan ``.''
  \item Validasi password: minimal 8 karakter, mengandung huruf dan angka
  \item Konfirmasi password: harus sama dengan password
  \item Tampilkan pesan error di \code{<span>} di bawah field menggunakan \code{textContent}
  \item Pada \code{submit}: periksa ulang semua field; \code{preventDefault()} jika gagal
\end{enumerate}
\end{contoh}

\begin{webcode}[language=JavaScript, caption={Validasi form real-time}]
const emailInput = document.querySelector("#email");
const errorEl = document.querySelector("#emailError");

emailInput.addEventListener("input", () => {
  const val = emailInput.value;
  if (!val.includes("@") || !val.includes(".")) {
    errorEl.textContent = "Format email tidak valid";
    emailInput.classList.add("error");
  } else {
    errorEl.textContent = "";
    emailInput.classList.remove("error");
  }
});
\end{webcode}

\begin{catatan}
\textbf{Validasi Client $\neq$ Validasi Server.} Validasi di browser meningkatkan UX tetapi \textbf{tidak} menggantikan validasi server. Pengguna dapat menonaktifkan JavaScript atau memodifikasi request. Validasi server wajib untuk keamanan; validasi client untuk kenyamanan. HTML5 juga menyediakan validasi bawaan (\code{required}, \code{type="email"}, \code{pattern}) yang dapat menjadi layer pertama \cite{mdnAccessibility}.
\end{catatan}

